% !TEX program = xelatex
% Independent consumer-side specification. Does not replace model design files.
\documentclass[10pt,UTF8,fontset=fandol]{ctexart}
\usepackage[a4paper,margin=22mm,headheight=15pt]{geometry}
\usepackage{amsmath,amssymb,booktabs,tabularx,array,enumitem,xcolor,fancyhdr}
\usepackage[unicode,colorlinks=true]{hyperref}
\definecolor{ink}{RGB}{24,53,77}
\definecolor{teal}{RGB}{26,120,113}
\hypersetup{linkcolor=ink,urlcolor=teal,pdftitle={TabU 预训练课程包：独立设计规范},pdfauthor={TabU Project}}
\ctexset{section={format=\Large\bfseries\color{ink}},subsection={format=\normalsize\bfseries\color{ink}}}
\setlength{\parskip}{0.25em}
\setlength{\emergencystretch}{3em}
\setlist{leftmargin=1.7em,itemsep=0.15em,topsep=0.2em}
\renewcommand{\arraystretch}{1.2}
\newcolumntype{Y}{>{\raggedright\arraybackslash}X}
\newcommand{\code}[1]{\texttt{\detokenize{#1}}}
\newcommand{\V}{\mathcal V}
\newcommand{\Q}{\mathcal Q}
\newcommand{\Obs}{\mathcal O}
\pagestyle{fancy}\fancyhf{}
\fancyhead[L]{\small TabU / Pretraining Curriculum}
\fancyhead[R]{\small 消费侧合同 / 2026-09-19}
\fancyfoot[C]{\thepage}
\begin{document}
{\LARGE\bfseries\color{ink} TabU 预训练课程包}\par
{\large 独立设计规范：冻结数据、可调用 episode 与证据阶梯}\par
\smallskip
\textbf{状态：}消费侧最小合同设计。本文规定数据与训练问题怎样登记、拆解和比较；
不报告新的训练效果，不把历史 TAR 成绩转写为 Gen-5.x 成绩。
现有可执行实现与本轮新增范围在第~\ref{sec:implementation} 节分别列明。

\section{三层设计与当前参考探针}
课程包应让不同 TabU 模型消费相同冻结数据、split、mask 与评估问题，独立替换编码、
backbone 或 readout。它是 \code{tabu-lab/src/tabu_lab/curriculum} 下的\textbf{消费侧模块}，
不新建仓库，不复制数据生成器。TFM-Data 继续拥有来源、world、纯数据供给与划分规则；
TabU 拥有模型适配、训练、checkpoint 与模型结果。核心合同不导入 Torch 或具体模型。

\textbf{A：不变语义。}来源与身份可追溯，模型可见证据和 scorer 真值分离；划分、比较条件、
初始化与实际曝光明确；结论不超出证据。\par
\textbf{B：可修订经验默认。}四阶段、8/120/120/12 表数、指标、mask 比例、预算和停止阈值
来自当前问题与历史诊断，可以按新证据修订；应记录解析后的选择、假设和变化理由。\par
\textbf{C：可替换部件。}生成器、数据来源、codec、模型、readout、trainer 与 scheduler
通过显式合同组合。替换部件不应暗中改变 A 层，也不要求永远沿用 B 层。

A 层的严格性服务于证据边界，而不是把研究过程写成机械流程。B 层除了可计量的默认值，
也容纳暂时不能阈值化的研究判断；它们可以先作为带有条件和出处的观察保留，等待更多证据
帮助澄清，而不必仓促伪装成统一规则。

\paragraph{当前课程的朴素主线。}
当前的方向先问一个很具体的问题：模型能否把一张训练表真正学会；随后才观察它能否在多张表上共同学习，
以及规模增大、持续学习中的旧表保持和真实数据适配会带来什么新的困难。只有这些层面各自有相应证据时，
未见表泛化才成为可以认真讨论的问题。这个主线用于组织问题和帮助定位失败，\textbf{不是}预先写死的执行
顺序、表数门槛或自动升级规则；各探针如何组合、何时扩展、哪些现象值得长期记录，仍应随着新的证据保留为
可修订的判断空间。

\begin{center}\small
\begin{tabularx}{\linewidth}{@{}p{24mm}YY@{}}\toprule
参考探针 & 当前默认回答的问题 & 必须保留的界限 \\\midrule
\code{fit8} & 八张合成表上的核心机制能否拟合 & 独立八模型与共享单模型须分开登记 \\
\code{multi120} & 单组参数能否共同拟合 120 张合成表 & 固定训练表拟合不等于未见 world 泛化 \\
\code{continual120} & 承接明确父权重后，能否学习严格分离的新 120 表 & 同时报告新表适应和旧表保留；继续旧表不算本阶段 \\
\code{real12} & 十二张真实表的训练池拟合，以及冻结保留行预测 & 拟合与 scorer-only reserved $R^2$ 独立出具回执 \\\bottomrule
\end{tabularx}
\end{center}

这些是\textbf{可组合的证据探针，不是强制线性 pipeline 或自动 checkpoint 链}。
可单独选择、改变顺序、扩展阶段或建立 factorized ablation。每次指定初始化来源：随机初始化、
weights-only 父权重或严格恢复。fit8 成功并不要求 multi120 必须承接其权重；历史 old120
从随机初始化开始，不能补写为连续课程。若声称新表持续学习，则须绑定实际父 checkpoint
及数据经历；不以父阶段的名字代替证据。旧表 replay 是单独的可调因素。

\subsection{通过标准必须在实验前明确}
每个探针登记问题、表集、预算、指标与阈值来源。manifest 被解析或校验通过不等于质量通过。
未设质量阈值时，运行完成只记
\code{completed}，不能记能力通过。训练池拟合指标可以回答容量诊断，不能挑选最终确认结果。
对于已知、可比较的问题，是否进入上一级应由已有证据与预先说明的判断依据决定，不因更新次数
耗尽自动升级。出现未预见的现象时，可先将其记为异常、未决问题或下一轮假设；不应为了填满
通过/失败栏而临时重写解释或阈值。

continual120 至少并列报告新表拟合变化与旧表退化。若记旧表损失为 $L_{\rm old}$，则
\[
\Delta L_{\rm old}=L_{\rm old}(\theta_{\rm after})-L_{\rm old}(\theta_{\rm before}).
\]
正值表示该固定诊断上的退化；其本身不能证明所有旧任务发生遗忘。换新数据时同时换优化器，
无法仅将改善归因于课程。课程收益需另设曝光与预算匹配的 scratch、joint 或顺序对照。

\newpage
\section{冻结 catalog：身份先于表名}
课程 manifest 引用原位冻结文件与已有 provenance，不搬运数据，不改写旧 manifest。
一个目录名、cohort 标签或不同 seed 都不足以证明数据独立。

\begin{center}\small
\begin{tabularx}{\linewidth}{@{}p{34mm}Y@{}}\toprule
身份层 & 最小记录与含义 \\\midrule
版本与来源 & catalog schema/version、dataset ID、来源 kind、原 manifest 路径及 digest；生成器版本或真实 upstream snapshot；缺失项显式声明 \\
冻结表 & 文件 URI/相对路径、原字节 SHA-256、表内容 fingerprint、typed schema；路径搬迁不改变数据身份 \\
World / snapshot & 合成 world 的已有身份或真实来源快照；生成请求与 seed 的出处；不能从文件名制造 world 身份 \\
表内 split & train/validation/reserved 原始行地址、split 规则与 seed（若有）、split digest；固定地址优先于重新随机切分 \\
消费阶段 & cohort、使用用途、明确表列表、预期表数与父阶段；已观察 reserved 标记为 retrospective \\\bottomrule
\end{tabularx}
\end{center}

\subsection{三种不同的隔离}
对于一张表的原始行地址集 $R$，要求
\[
R_{\rm train}\;\dot\cup\;R_{\rm validation}\;\dot\cup\;R_{\rm reserved}=R.
\]
地址唯一、合法、互斥且完整覆盖。此条件只证明\emph{表内行隔离}；不证明 world 或来源独立。
自然缺失与人工隐藏另记，缺失值不得变成有监督目标。

当一个新表探针声明严格分离时，必须相对于声明的前序训练表检查身份集合：表 ID 别名、原字节、
规范化表内容和可获得的 world/snapshot 身份。首版内容哈希仅使用 values：保留列顺序，
统一整数/浮点表示，对完整行排序后计算；排除 schema、split、格式空白和表名。
因此相同 values 即使改变 schema 也会被保守拒绝；不声称检测列置换或任意变换。
原始行顺序、类型和 split 地址另由来源字节与 split digest 绑定。
相同 world 的新行应登记为同世界再采样，不能自动称为未见世界。

分离结论必须注明证据层：\textbf{table-content-disjoint} 与
\textbf{world-disjoint} 是两种声明。严格声明所需身份缺失时失败并报告缺项，
不凭种子不同补足证明。旧数据缺少 world 身份时可保留表级诊断用途，但不能获得世界级独立标签。
同一 cohort 内的重复也须显式去重或声明权重，避免别名让表数和训练曝光虚增。

\subsection{版本与重放}
manifest 的解析结果、数据与 split fingerprint、episode 配方、独立 seed streams
共同确定课程身份。生成器、来源或 split 变化均生成新身份。原始字节哈希绑定版本，
内容哈希辅助识别别名；二者不互相代替，也不单独证明 upstream 来源真实。

校验器输出可读的\emph{缺失、冲突、已验证层次}，不重新生成一批相似数据填补历史缺口。
读取 legacy 表需保留到原始行列的映射；TFM-Data 的数据 schema、历史 TAR typed JSON
与模型输入对象分属不同合同，不因字段相近而隐式互换。

\newpage
\section{可调用 episode 与 scorer 真值边界}
EpisodeFactory 接收冻结表引用、partition/purpose、episode index、配方与随机种子，
返回分离的模型证据、目标地址、scorer 真值和无真值的 provenance receipt。
核心接口描述语义；具体张量化与模型对象转换放在可选适配器。

令 $\Obs$ 为自然观测 cell，$\Q\subseteq\Obs$ 为人工隐藏 Query，$\V=\Obs\setminus\Q$。
当前全观测恢复默认请求 $T=\Obs$；目标集合与 state 权重属于可修改配方。
V4 的 Query-only 与 V5 的 all-observed 目标不同，不能默认为相同任务：
\[
E_{\rm in}=\bigl(\mathrm{schema},\V,\Q,\{x_{ra}:(r,a)\in\V\}\bigr),
\qquad T=\Obs,\qquad Y^\star=\{x_{ra}:(r,a)\in T\}.
\]
Query 和 Null 的 payload 在进入模型前物理擦除；隐藏值仅在 \code{TruthSidecar} 中保留。
Forward、codec、支持选择与 readout 只接收 $E_{\rm in}$；真值在 scorer 与预测首次汇合。
\[
\widehat Y=F_\theta(E_{\rm in},T),\qquad
\mathcal L=\mathrm{score}(\widehat Y,Y^\star).
\]
固定输入证据后替换 sidecar 真值，forward 必须不变。\code{detach()} 只处理梯度，不能替代隔离。

\begin{center}\small
\begin{tabularx}{\linewidth}{@{}p{34mm}Y@{}}\toprule
接口对象 & 内容与允许的消费者 \\\midrule
Episode request & 表与 split 身份、用途、index、recipe、seed namespace；由课程调度生成 \\
模型证据与目标 & typed 可见值、visible/query/observed masks、目标地址；仅此部分给模型 \\
Scorer truth & 原始观测真值与目标状态；只给 loss/evaluator，不放入一般 receipt \\
Episode receipt & 数据/split/recipe 身份、实际行列映射、Query 地址、支持计数、有效 seed 与拒绝原因；不含隐藏值 \\\bottomrule
\end{tabularx}
\end{center}

\subsection{独立随机性与可复现性}
把模型初始化、表顺序、窗口、mask、离散编码和评价分成独立 streams。
生成器与 split 的 seed 属于数据 provenance，不挤进训练 RNG。一次事件可寻址为
\[
\xi_k=H(s_k,\mathrm{stage},\mathrm{table},\mathrm{partition},\mathrm{index},\mathrm{purpose}).
\]
命名固定 probe 使用固定 namespace，不随训练 update 重抽。receipt 记录实际地址与种子，
窗口可重放不等于全行覆盖。恢复接续绝对 cursor，不重新从首个 episode 抽样。

\subsection{共享任务语义与模型专属数值编码}
每列支持数、类别覆盖、Query 比例和目标集合是任务规则，必须显式记录。
若配方声明精确 Query 预算，则本适配器在无法满足时整次 episode 失败；其他接受、拒绝或
再采样政策可替换，但须显式记录，不能静默缩小 loss 或删除难例。
numeric tail guard 是改变问题的配方因素，默认关闭；reserved evaluation 禁止据标签筛难例。
median/half-IQR 等统计只能由 episode-visible 证据计算，并在该 episode 的编码、scorer、
decode 中一致使用。编码维度、基函数与 LL/NW 选择归模型配置，不固化到课程核心。

\newpage
\section{真实保留评估、因素分解与回执}
\subsection{real12：训练拟合与 reserved \texorpdfstring{$R^2$}{R2} 分开}
训练仅从 train 池构造 restoration episode。冻结评价时隐藏 reserved 的全部目标列，
目标真值只交 scorer。若 reserved 的其他特征可见，必须标记 transductive；
单条 Query 与联合 Query 是不同协议。已反复查看或用于选模型的历史 reserved 只能作为
retrospective 诊断；需要新的、明确未接触确认集才能重获最终确认含义。

对表 $t$ 的固定数值目标地址 $Q_t$，
\[
R_t^2=1-\frac{\sum_{i\in Q_t}(y_i-\hat y_i)^2}
{\sum_{i\in Q_t}(y_i-\bar y_t)^2},
\qquad \bar y_t=\frac{1}{|Q_t|}\sum_{i\in Q_t}y_i.
\]
保留逐表 $R_t^2$、目标数、分母与失败状态；分母为零时报告 undefined。
宏平均必须同时给出有效表数和缺失原因，不能静默跳过失败表后称为完整 real12。
分类表另报适当指标，不把编码 MSE 替代原值 $R^2$。

\subsection{可替换因素的边界}
\begin{center}\small
\begin{tabularx}{\linewidth}{@{}p{28mm}YY@{}}\toprule
因素 & 课程/数据合同 & 模型或执行合同 \\\midrule
数据供给 & 表集、来源版本、split、world 关系 & 不控制生成器或重新抽取数据 \\
任务 & mask、窗口、目标、支持与类别覆盖、probe & 编码、backbone、LL/NW、损失实现 \\
调度 & cohort 权重、顺序、replay、stage 选择 & optimizer、学习率、精度、设备、batch \\
证据 & 前后固定探针、隔离范围、预注册停止规则 & 父 checkpoint、实际更新与耗时、运行失败 \\\bottomrule
\end{tabularx}
\end{center}
完整实验组合两侧配置并绑定身份；修改任意影响执行的因素形成新比较身份。
多因素同时变化可作为实用方案，但不能归因于某个单一因素。

\subsection{最小结果语义与停止推进规则}
一次 receipt 至少含 schema/version、课程与数据身份、model/source/config 身份、
父 checkpoint、stage/question、实际曝光、预算、状态、逐表指标、原始产物引用与 digest。
\code{planned}、\code{completed}、\code{stopped}、\code{failed} 和质量 verdict 分开；
实现检查通过、短 fit、保留评价与可泛化能力分别陈述。

训练停在预注册 updates/time 上限、数值失败、数据身份漂移或阶段质量停止条件。
失败保留最后有效 checkpoint 与 terminal receipt；不自动延长预算或改阈值寻找通过。停止一条
路线并不否定其所有价值：负结果、异常模式和暂时无法归因的变化也应与其上下文一同保留，
成为后续研究判断的材料，而非被当作噪声抹去。
训练池 fit gate 只能认证该层拟合；validation 可以参与选型；final reserved 不能驱动调参。
严格 resume 恢复 optimizer/RNG/cursor；weights-only 初始化另建身份，不能伪装成 resume。

\newpage
\section{既有执行器与本轮最小实现边界}\label{sec:implementation}
\subsection{先复用已存在的 V5.3 工作}
当前 \code{tabu-lab} 已有独立 \code{curriculum_v53}，由提交 \code{3266d1b}
引入并在 \code{6293e95} 合并。它提供数据 SHA 校验、严格行划分、六路 seed、
manifest 调度、固定 probe、FP64 本地 runner、原子 checkpoint 与严格恢复。
其 \code{data/protocol/runner/evaluation} 直接依赖 V5.3；这是一条模型专属执行路径，
不能直接称为跨模型课程包。其源码身份和历史恢复语义保持原样。

现有实现保护训练行与 scorer truth 的边界，但 cohort 名称不证明新表独立：
表条目目前只检查 ID 唯一，允许相同数据字节以不同别名加入 old/new。
现有调度测试也刻意复用同一微型表来验证 120 表轮换。这是调度测试，不能拿来证明数据分离。
本轮最小合同优先补齐\textbf{冻结 catalog、跨集合分离、阶段选择及公共 request/receipt}，
不再制造第二套训练器。

\subsection{本轮最小公开接口}
\begin{center}\small
\begin{tabularx}{\linewidth}{@{}p{44mm}Y@{}}\toprule
接口 & 边界 \\\midrule
\code{bind_table} & 校验预期文件 SHA，绑定表、cohort、kind 与 provenance，返回冻结 \code{BoundTable} \\
\code{bind_legacy_corpus} & 引用已有 corpus manifest，保留其 SHA、record、seed 与可用 world 身份 \\
\code{require_disjoint} & \code{level=table} 校验 bytes/content 及已知 world 冲突；\code{level=world} 另要求身份完整 \\
\code{select_stage} & 通用阶段选择：任意 name/question，可选 expected count、kind、previous、separation level 与 reserved 要求 \\
\code{reference_stage} & 提供 fit8/multi120/continual120/real12 的当前参考默认；允许显式覆盖并记录解析后的选择 \\
\code{EpisodeRequest} & 配方、index、partition/purpose 与 masks/codes/windows/evaluation 种子；不混入模型配置 \\
\code{V53EpisodeFactory} & 可选 lazy bridge，复用现有 \code{load_table/build_episode}，返回 inputs/request/truth/audit \\\bottomrule
\end{tabularx}
\end{center}
上述接口属于 \code{tabu_lab.curriculum}。其中 catalog/阶段对象是纯 metadata，
EpisodeFactory 为调用合同；V5.3 依赖只在可选 bridge 进入。API 接收当前已有 provenance，
不会因字段被填写就证明 upstream 来源正确。

\subsection{实现约束与后续验收}
核心不加载模型，不执行 optimizer，不生成数据，不暴露 reserved 真值。
适配器显式复用已有 mask/episode 实现；不修改原 runner 的源码哈希或恢复协议。
模块验收覆盖别名/内容/world 冲突、split 重叠、seed 重放、缺失 provenance、
显式数量/用途、自定义阶段与默认覆盖、隐藏值隔离，以及核心 import 不加载 Torch/模型。
这些检查不构成训练成功。TFM-Data 持续供给、世界级训练/确认接入、GPU/大表与分布式执行
需独立验证；本轮不训练、不移动数据、不改历史记录。

结果 schema、自动质量 gate 和停止调度在本文中是建议合同，\textbf{本轮没有另写执行器}。
此 metadata wedge 不证明任意第三方 factory 都正确实现隔离；具体 adapter 仍需 truth
替换不影响 forward 的测试。现有 V5.3 adapter 不能直接装载 legacy TAR checkpoint。
它支持的 recipe 种类只是本轮实现范围，不是对未来课程或其他模型的通用限制。
首版可解析认证的 world 身份来自既有 legacy corpus record；其他来源的 provenance
仍属引用与声明，不能仅靠自填不同 ID 通过严格 world 分离。

\newpage
\section{历史证据带来的具体设计选择}
下表列出已存在的原始回执所支持的有限结论；不用于构造 V4 与 Gen-5.x 的等预算排名。
原始 archive/research-proposals 路径均相对 TabU 项目根，\code{docs/} 与 \code{src/}
路径则相对嵌套 \code{tabu-lab}。数据和 checkpoint 均原位引用。

\begin{center}\small
\begin{tabularx}{\linewidth}{@{}p{23mm}YY@{}}\toprule
历史诊断 & 记录中已观察到的结果 & 对课程合同的直接要求 \\\midrule
V4 八表独立拟合 & 八个模型各 1,024 updates；76 行训练池；数值 NMSE 约 0.00145--0.01582，XOR 100\% & 不能把八个独立模型叫作共享 fit8；不预填任意近零阈值 \\
V4 old120 & train768 为 768 rounds / 92,160 updates；固定训练行；median NMSE 约 0.03824 & rounds 和 updates 并存；固定训练拟合不能充当 reserved 结果 \\
V4 新表无 replay & 3,840 updates 后新表 median accuracy 60.48\% 至 90.99\%；旧表 99.26\% 至 80.88\% & 新表适应与旧表保留必须并列；仅新表改善不算完整持续学习结论 \\
历史 real12 & legacy TAR reserved macro $R^2\approx0.76006$；后续 Restoration 全宽 continuation 约 0.65101 & 历史 reserved 已被观察，只能 retrospective；模型、目标和训练历程并不匹配 \\\bottomrule
\end{tabularx}
\end{center}

\begingroup\raggedright
八表记录来自 \path{archive/tar-unified-synthetic-nearzero-20260907/REPORT.md}；
old120 来自 \path{archive/tar-diverse-120-fit-20260907/train768.REPORT.md}；
新表与保留变化来自 \path{archive/tar-new120-muon-20260908/REPORT.md}。
该新表实验还记录无 values、retry seed window 与可用机制身份重叠的 corpus audit，
支持把\emph{实际身份审计}设为 continual120 的必需输入，而非依赖不同表名。
\par\endgroup

本轮消费侧 metadata 审计核对 old/new 共 240 张表的文件 SHA 并通过表内容分离检查；
其中 96 张可由当前解析器认证 \code{world_hash}，144 张未取得这一层证明，严格 world
分离因此不通过。这只说明当前解析范围，不能推断原始来源没有其他机制身份。

real12 数字来自 \path{research-proposals/restoration-real12-rescue-20260917/REPORT.md}
的 2026-09-18 effective-direction continuation terminal result；该节同时记录
train macro $R^2$ 为 0.7330293、legacy 参照为 0.7541652。训练改善没有同步提高 reserved，
所以两类结果必须分开报告。现有证据不支持只凭数字差异归因于某一编码或 readout 改动。

保留原实验名称、source revision、命令、数据与 checkpoint 身份，并登记证据用途。
新课程配置只引用旧数据与回执；历史 8/120/12 的结果不能被追溯写成一个已证明有效的统一课程。
Gen-5.x 的编码与预测结构变化归模型设计记录，课程核心不借版本号改变 restoration 任务。
当前 V5.3 新训练效果尚未在本规范中核验，不从本地 runner 资格推断其训练成功。

\subsection{可核查的当前来源}
\begin{itemize}\small
\item TabU 项目 \code{README.md} 第 7、127 行及 TFM-Data \code{data-contract.md}：数据供给与模型消费归属。
\item \code{docs/tutorials/v53-curriculum.md}：现有 runner 的 probe、seed、初始化与恢复语义。
\item \code{docs/reviews/v53-curriculum-20260919/README.md}：已记录的本地执行资格与限制；本文未重跑训练。
\item \code{src/tabu_lab/curriculum_v53/data.py}、\code{protocol.py}；
      \code{tests/unit/curriculum_v53/test_protocol.py}：当前数据/split 校验与调度测试。
\item \code{src/tabu_lab/models/restoration/contracts.py}：物理擦除隐藏值、完整观测目标与 TruthSidecar。
\end{itemize}
\end{document}
